\section{Opis architektury}
\label{sec:architektura}

System składa się z aplikacji klienckiej działającej w przeglądarce (PWA),
serwera REST API, relacyjnej bazy danych oraz magazynu obiektów na pliki
dokumentów. Wszystkie komponenty działają w kontenerach Docker i są
uruchamiane wspólnie przez Docker Compose.

Rysunek~\ref{fig:kontenery} przedstawia kontenery systemu w notacji modelu C4.
Użytkownik łączy się wyłącznie z kontenerem \texttt{frontend}:
serwer nginx udostępnia zbudowaną aplikację React, a żądania na ścieżki
\texttt{/api/*} przekazuje do kontenera \texttt{backend}. Przeglądarka rozmawia
więc z jednym adresem (ta sama domena dla aplikacji i API), dzięki czemu nie
trzeba konfigurować mechanizmu CORS. Backend zapisuje dane w bazie PostgreSQL
i zaszyfrowane pliki w magazynie MinIO, zgodnym z interfejsem Amazon S3.
Prometheus okresowo pobiera metryki backendu, a Grafana prezentuje je na
dashboardach.

\begin{figure}[htbp]
\centering
\includegraphics[width=0.85\textwidth]{diagramy/kontenery.png}
\caption{Diagram kontenerów systemu (model C4)}
\label{fig:kontenery}
\zrodlo{opracowanie własne}
\end{figure}

Kontenery są rozdzielone na trzy sieci Dockera według zasady najmniejszych
uprawnień: \texttt{app} (nginx i backend), \texttt{data} (backend, PostgreSQL,
MinIO) oraz \texttt{monitoring} (backend, Prometheus, Grafana). Frontend nie ma
bezpośredniego dostępu do bazy danych ani do magazynu plików, a jedynym
kontenerem obecnym we wszystkich sieciach jest backend. Szczegóły konfiguracji
opisano w podrozdziale~\ref{sec:docker}.

\subsection{Wybór technologii}

W całym projekcie użyto języka TypeScript. Frontend i backend
korzystają dzięki temu z tych samych typów i schematów walidacji, umieszczonych
we wspólnym pakiecie. Zestawienie technologii zawiera
tabela~\ref{tab:technologie}.

\begin{table}[htbp]
\centering
\caption{Technologie użyte w projekcie}
\label{tab:technologie}
\small
\begin{tabular}{|p{0.2\textwidth}|p{0.3\textwidth}|p{0.38\textwidth}|}
\hline
\textbf{Obszar} & \textbf{Technologia} & \textbf{Uzasadnienie} \\
\hline
Backend & Node.js 22, Express & Wskazany w wymaganiach, prosty serwer REST \\
\hline
Baza danych & PostgreSQL 16 & Wskazana w wymaganiach \\
\hline
ORM i migracje & Prisma & Schemat bazy w jednym pliku, migracje wersjonowane w repozytorium \\
\hline
Pliki & MinIO (S3) & Magazyn obiektów w osobnym kontenerze \\
\hline
Walidacja & Zod & Te same reguły na froncie i w API \\
\hline
Frontend & React 19, Vite & Szybkie środowisko deweloperskie, łatwa konfiguracja PWA \\
\hline
Uwierzytelnianie & JWT, argon2 & Tokeny z wymagań; aktualny standard skracania haseł \\
\hline
Szyfrowanie & AES-256-GCM (\texttt{node:crypto}) & Szyfrowanie kopertowe plików przed zapisem \\
\hline
Monitoring & Prometheus, Grafana & Metryki, stan zdrowia, alerty, dashboardy \\
\hline
Konteneryzacja & Docker, Docker Compose & Uruchomienie całego systemu jednym poleceniem \\
\hline
Testy & Vitest, Supertest, Playwright & Testy funkcjonalne API i interfejsu, testy jednostkowe \\
\hline
\end{tabular}
\zrodlo{opracowanie własne}
\end{table}

\subsection{Struktura repozytorium}

Kod jest zorganizowany jako monorepo oparte na mechanizmie \emph{workspaces}
menedżera npm. Jedno polecenie \texttt{npm install} instaluje zależności
wszystkich pakietów, a wspólne skrypty (\texttt{build}, \texttt{lint},
\texttt{test}) uruchamiane są z katalogu głównego. Podział katalogów
przedstawia tabela~\ref{tab:repozytorium}.

\begin{table}[htbp]
\centering
\caption{Katalogi repozytorium}
\label{tab:repozytorium}
\small
\begin{tabular}{|p{0.22\textwidth}|p{0.66\textwidth}|}
\hline
\textbf{Katalog} & \textbf{Zawartość} \\
\hline
\texttt{backend/} & Serwer REST API (Express) \\
\hline
\texttt{frontend/} & Aplikacja kliencka React (PWA) \\
\hline
\texttt{shared/} & Schematy walidacji, typy i wyliczenia wspólne dla frontu i backendu \\
\hline
\texttt{e2e/} & Funkcjonalne testy interfejsu (Playwright) \\
\hline
\texttt{monitoring/} & Konfiguracja Prometheusa, Grafany i alertów \\
\hline
\texttt{docs/} & Niniejsze sprawozdanie (\LaTeX) i diagramy \\
\hline
\texttt{scripts/} & Skrypty pomocnicze, m.in. budowanie sprawozdania w Dockerze \\
\hline
\end{tabular}
\zrodlo{opracowanie własne}
\end{table}

Pakiet \texttt{shared} jest kompilowany jako pierwszy, ponieważ importują go
backend i frontend. Jakość kodu pilnują ESLint (reguły dla TypeScriptu
i Reacta) oraz Prettier (jednolite formatowanie); oba narzędzia uruchamia
polecenie \texttt{npm run lint}. Wspólne ustawienia kompilatora TypeScript
(tryb \texttt{strict}) znajdują się w pliku \texttt{tsconfig.base.json},
rozszerzanym przez każdy pakiet.
